\section{微扰论中的 \qcd{} 梯度流}\label{sec:pert}

为完整起见，本节汇集 \citeres{Luscher:2010iy,Luscher:2011bx} 原始推导的
主要步骤。


\subsection{梯度流的定义}

下面我们在 $D$ 维欧氏时空中工作，$D=4-2\ep$。\gff\ 把常规
（regular）\footnote{我们用``流动的''（flowed）与``常规的''（regular）
\qcd\ 来区分在 $t>0$ 与在 $t=0$ 定义的量。}\qcd\ 的胶子场与夸克场
$A^a_\mu(x)$、$\psi_\alpha^i(x)$ 延拓为 $(D+1)$ 维场
$B^a_\mu(t,x)$、$\chi_\alpha^i(t,x)$，其边界条件为
\begin{equation}
  \begin{split}
    B_\mu^a (t=0,x) = A_\mu^a (x)\,,\qquad \chi_\alpha^i (t=0,x)=
    \psi_\alpha^i(x)
    \label{eq:bound}
  \end{split}
\end{equation}
流方程为\,\cite{Luscher:2010iy,Luscher:2013cpa}
\begin{equation}
  \begin{split}
     \partial_t B_\mu^a &= \mathcal{D}^{ab}_\nu G_{\nu\mu}^b + \kappa
     \mathcal{D}^{ab}_\mu \partial_\nu B_\nu^b\,,\\ \partial_t \chi &= \Delta \chi
     - \kappa \partial_\mu B_\mu^a T^a \chi\,,\\ \partial_t \bar
     \chi &= \bar \chi \overleftarrow \Delta + \kappa \bar \chi
     \partial_\mu B_\mu^a T^a\,,
    \label{eq:flow}
  \end{split}
\end{equation}
其中``流时间'' $t$ 是质量维度为负二的参数，$\kappa$
是一个规范参数，它从物理可观测量中消失（见下文）。

$(D+1)$ 维场强张量定义为
\begin{align}
  G_{\mu\nu}^a = \partial_\mu B_\nu^a - \partial_\nu B_\mu^a + f^{abc}
  B_\mu^b B_\nu^c\,,
\end{align}
伴随表示的协变导数为
\begin{align}
  \mathcal{D}_\mu^{ab} = \delta^{ab} \partial_\mu - f^{abc} B_\mu^c\,,
\end{align}
而
\begin{equation}
  \begin{split}
    \Delta = \mathcal{D}^\mathrm{F}_\mu \mathcal{D}^\mathrm{F}_\mu
    \,,\qquad
    \overleftarrow{\Delta} = \overleftarrow{\mathcal{D}}^\mathrm{F}_\mu
    \overleftarrow{\mathcal{D}}^\mathrm{F}_\mu \,,
  \end{split}
\end{equation}
其中基本表示的协变导数为
\begin{equation}
  \mathcal{D}^\mathrm{F}_\mu = \partial_\mu + B^a_\mu T^a \,,\qquad
  \overleftarrow{\mathcal{D}}^\mathrm{F}_\mu =
    \overleftarrow{\partial}_\mu - B^a_\mu T^a \,.
\end{equation}
按照惯例，伴随表示的色指标记为 $a,b,c,\ldots$，而 $\mu,\nu,\rho,\ldots$
为 $D$ 维洛伦兹指标。基本表示的色指标记为 $i,j,k,\ldots$，但除非清晰所需，
全文将其省略；自旋指标 $\alpha,\beta,\gamma,\ldots$ 亦同。对称性生成元
$T^a$ 取基本表示。它们满足对易关系
\begin{equation}
  \begin{split}
    [T^a,T^b] = f^{abc}T^c\,,
  \end{split}
\end{equation}
其中 $f^{abc}$ 为结构常数，迹归一化为
\begin{equation}
  \mathrm{Tr}(T^a T^b) = - \TR \delta^{ab} \, ,
\end{equation}
$\TR > 0$。

\noeqn{eq:flow} 中 $\kappa$ 的不同取值对应如下形式的规范变换：
\begin{align}
  \chi \rightarrow \chi^\prime = \Lambda \chi \quad \quad \text{及}
  \quad \quad B_\mu^a T^a \rightarrow B_\mu^{a\prime} T^a= \Lambda
  B_\mu^a T^a \Lambda^{-1} + \Lambda \partial_\mu \Lambda^{-1}\,,
\end{align}
其中
\begin{align}
  \Lambda(t,x)=e^{-\int_0^t \mathrm{d}s \, \kappa \partial_\mu B_\mu^a(s,x) T^a}\,.
\end{align}
当然，所有可观测量都与规范参数 $\kappa$ 无关\,\cite{Luscher:2010iy}。
在微扰计算中，通常取 $\kappa = 1$ 最为方便。

通过定义
\begin{align}
  \mathcal{L} = \mathcal{L}_\mathrm{\qcd} +
  \mathcal{L}_\textrm{gauge-fixing} + \mathcal{L}_\mathrm{ghost} +
  \mathcal{L}_{B} + \mathcal{L}_{\chi}.
  \label{eq:Lagrangian_complete}
\end{align}
可以把流方程 \noeqn{eq:flow} 纳入拉格朗日形式体系。前三项构成常规的
Yang-Mills 拉格朗日量，并在基本表示中加入费米子（夸克）。引入指标 $f$
以区分质量为 $m_f$ 的不同夸克味，经典部分、规范固定部分与 Faddeev-Popov
鬼部分分别由下式给出：
\begin{equation}
  \begin{split}
  \mathcal{L}_\mathrm{\qcd} &= \frac{1}{4g^2} F_{\mu \nu}^a F_{\mu
    \nu}^a + \sum_{f=1}^{\NF}\bar{\psi}_f ( \slashed{D}^\mathrm{F}+m_f)
  \psi_f\,,\\ \mathcal{L}_\textrm{gauge-fixing}&= \frac{1}{2 g^2 \xi}
  (\partial_\mu A_\mu^a)^2\,,\\ \mathcal{L}_\mathrm{ghost} &=
  \frac{1}{g^2}\partial_\mu \bar{c}^a D_\mu^{ab} c^b\,,
  \label{eq:lags}
  \end{split}
\end{equation}
其中
\begin{equation}
  \begin{split}
    F_{\mu\nu}^a = \partial_\mu A^a_\nu - \partial_\nu A^a_\mu + f^{abc}A_\mu^bA_\nu^c
    %% \label{eq:}
  \end{split}
\end{equation}
是常规场强张量，而
\begin{equation}
  D_\mu^\mathrm{F} = \partial_\mu + A_\mu^a T^a \,,\qquad D_\mu^{ab} = \delta^{ab} \partial_\mu - f^{abc} A_\mu^c
\end{equation}
分别是基本表示与伴随表示的常规协变导数；$g$ 为规范耦合，$\xi$ 为
\qcd\ 规范参数，$\NF$ 为不同夸克味的数目。流方程通过引入拉格朗日乘子场
\begin{equation}
  \begin{split}
    L_\mu^a(t,x)\qquad\text{与}\qquad
    \lambda_f(t,x)\,,\ \bar\lambda_f(t,x)\,,
  \end{split}
\end{equation}
纳入体系；除质量维度分别为 $3$ 和 $5/2$ 外，它们分别携带与流动后的胶子场和
夸克/反夸克场 $B_\mu^a$、$\chi,\bar\chi$ 相同的量子数。由下式
\begin{equation}
  \begin{split}
  \mathcal{L}_{B} &= -2 \int_0^\infty \mathrm{d}t \, \textrm{Tr} \left[
    L_\mu^a T^a \left(\partial_t B_\mu^b T^b -\mathcal{D}_\nu^{bc}
    G_{\nu \mu}^c T^b - \kappa \mathcal{D}_\mu^{bc} \partial_\nu B_\nu^c
    T^b \right) \right]\,,\\ \mathcal{L}_{\chi} & =
  \sum_{f=1}^{\NF}\int_0^\infty \mathrm{d}t \, \Big( \bar{\lambda}_f
  \left(\partial_t - \Delta + \kappa \left(\partial_\mu B_\mu^a\right)
  T^a \right) \chi_f \\&\hspace*{10em}
  + \bar{\chi}_f \left(\overleftarrow{\partial_t} -
  \overleftarrow{\Delta} - \kappa \left(\partial_\mu B_\mu^a\right) T^a
  \right) \lambda_f \Big)\,,
  \label{eq:Lagrangian_flow}
  \end{split}
\end{equation}
导出的 Euler-Lagrange 方程确实给出 \noeqn{eq:flow}\,\cite{Luscher:2011bx}。

注意，拉格朗日量 \noeqn{eq:Lagrangian_complete} 并不包含流动鬼场
$d^a(t,x)$ 与 $\bar{d}^a(t,x)$。它们与通常的 Faddeev-Popov 鬼 $c^a(x)$、
$\bar{c}^a(x)$ 一样因规范固定而产生，并满足初始条件
\begin{align}
  \left. d^a (t,x) \right|_{t=0} = c^a(x).
\end{align}
与常规规范场的鬼类似，只要没有外线 $d^a$ 或 $\bar{d}^a$（在含外胶子的振幅中
只考虑物理自由度即可避免），它们总是形成闭合圈。后面会看清楚：仅由流动场构
成的闭合圈为零，因此可以在拉格朗日量层面就把 $d^a$ 与 $\bar{d}^a$ 略去。


\subsection{流方程的微扰解}

让我们引入简记
\begin{equation}
  \begin{split}
    \int_p \equiv \int\frac{\dd^D p}{(2\pi)^D}\,,\qquad
    \int_x \equiv \int\dd^D x\,,
    %% \label{eq:}
  \end{split}
\end{equation}
积分变量处于坐标空间（$x,y,z,\ldots$）还是动量空间（$p,k,q,\ldots$）应由上
下文自明。

把流动规范场 $B_\mu^a(t,x)$ 的流方程 \noeqn{eq:flow} 分成线性与非线性两部分
是有帮助的：
\begin{equation}
  \begin{split}
  &\partial_t B_\mu^a = \partial_\nu \partial_\nu B_\mu^a + (\kappa - 1)
  \partial_\mu \partial_\nu B_\nu^a + R_\mu^a, \\ &R_\mu^a = 2 f^{abc}
  B_\nu^b \partial_\nu B_\mu^c - f^{abc} B_\nu^b \partial_\mu B_\nu^c +
  (\kappa - 1) f^{abc} B_\mu^b \partial_\nu B_\nu^c + f^{abe} f^{cde}
  B_\nu^b B_\nu^c B_\mu^d.
  \label{eq:flowsep}
  \end{split}
\end{equation}
线性方程可以通过引入积分核
\begin{equation}
  K_{\mu \nu} (t,x) = \int_p \frac{e^{\mathrm{i}px}}{p^2} \left( \left(
  \delta_{\mu \nu} p^2 - p_\mu p_\nu \right) e^{-tp^2} + p_\mu p_\nu \,
  e^{-\kappa tp^2} \right) \equiv \int_p e^{\ii px}\widetilde K_{\mu\nu}(t,p)\,,
  \label{eq:kmunu}
\end{equation}
求解，它满足
\begin{equation}
  \lim_{t \to 0} K_{\mu\nu}(t,x) = \delta_{\mu\nu} \delta^{(D)}(x).
\end{equation}
计入初始条件 \noeqn{eq:bound} 后，流方程的完整解即为
\begin{equation}
  B_\mu^a (t,x) = \int_y K_{\mu \nu}(t,x-y) A_{\nu}^a(y)+ \int_y
  \int_0^t \mathrm{d}s \, K_{\mu \nu}(t-s,x-y) R_\nu^a(s,y)\,,
\end{equation}
或在动量空间中，
\begin{equation}
  \widetilde{B}_\mu^a(t,p)= \int_x e^{-\mathrm{i}px} B_\mu^a(t,x)=
  \widetilde{K}_{\mu \nu}(t,p) \widetilde{A}_\nu^a(p)+ \int_0^t
  \mathrm{d}s \, \widetilde{K}_{\mu \nu}(t-s,p)
  \widetilde{R}_\nu^a(s,p).
  \label{eq:gluon_solution_momentum}
\end{equation}
把该解迭代地代入自身，可将 \noeqn{eq:flowsep} 非线性部分的傅里叶变换表示为
\begin{equation}
  \begin{split}
  \widetilde{R}_\mu^a(t,p) &=
  \int_{q,l,k} (2 \pi)^D \, \delta^{(D)}(p-q-l-k)\bigg[
   \delta^{(D)}(k)\cdot X_{2,\mu\nu\rho}^{abc}(q,l)\, \widetilde{B}_{\nu}^{b}(t,q)
      \widetilde{B}_{\rho}^{c}(t,l)\\&
  +
   X_{3,\mu\nu\rho\sigma}^{abcd}\, \widetilde{B}_{\nu}^{b}(t,q)
      \widetilde{B}_{\rho}^{c}(t,l)
      \widetilde{B}_{\sigma}^{d}(t,k)\bigg]\,,
    %% \label{eq:}
  \end{split}
\end{equation}
其中
\begin{equation}
  \begin{split}
  X_{2,\mu \nu \rho}^{abc}(q,l) &= -\mathrm{i} f^{a b c} \big[
  (l-q)_\mu \delta_{\nu \rho} + 2q{}_{\rho} \delta_{\mu \nu}- 2l{}_{\nu}
  \delta_{\mu \rho} + (\kappa - 1) ( q{}_{\nu} \delta_{\mu
    \rho} -l{}_{\rho} \delta_{\mu \nu}) \big]\,,\\
  X_{3,\mu \nu \rho \sigma}^{abcd} &= f^{abe}
  f^{cde}(\delta_{\mu \sigma} \delta_{\nu \rho} - \delta_{\mu \rho}
  \delta_{\nu \sigma}) + f^{ade} f^{bce}(\delta_{\mu \rho}
  \delta_{\nu \sigma}- \delta_{\mu \nu} \delta_{\rho \sigma})\\&\quad
  + f^{ace} f^{dbe}(\delta_{\mu \nu} \delta_{\rho \sigma}-
  \delta_{\mu \sigma} \delta_{\nu \rho}).
    %% \label{eq:}
  \end{split}
\end{equation}
$X_3$ 的结构与常规 \qcd\ 的四胶子顶点完全相同。稍后表述费曼规则时，$X_2$
与 $X_3$ 描述流动胶子场的三点与四点顶点。

流动夸克场的流方程 \noeqn{eq:flow} 同样可分解为线性与非线性部分来求解：
\begin{align}
  \partial_t \chi =\partial_\mu \partial_\mu \chi + \Delta^\prime \chi
  \quad \text{其中} \quad \Delta^\prime = (1-\kappa) \partial_\mu
  B_\mu^a T^a + 2 B_\mu^a T^a \partial_\mu + B_\mu^a B_\mu^b T^a T^b\,.
\end{align}
线性方程由积分核
\begin{equation}
  K(t,x)=\int_p e^{\mathrm{i}px} e^{-tp^2} \equiv \int_p e^{\ii
    px}\widetilde K(t,p)\,,
  \label{eq:k}
\end{equation}
求解，借助它可以写出完整解
\begin{equation}
  \chi(t,x)= \int_y K(t,x-y)\psi(y) + \int_y \int_0^t \mathrm{d}s \, K(t-s,x-y) \Delta^\prime \chi(s,y).
\end{equation}
此处及下文，除非清晰所需，我们略去味指标 $f$。傅里叶变换后的场
\begin{equation}
  \widetilde{\chi}(t,p)= \widetilde{K}(t,p)\widetilde{\psi}(p) + \int_0^t \mathrm{d}s \, \widetilde{K}(t-s,p)\widetilde{\Delta^\prime \chi}(s,p)
  \label{eq:quark_solution_momentum}
\end{equation}
的非线性部分可以表示为
\begin{equation}
  \begin{split}
    \widetilde{\Delta^\prime \chi}(t,p) =&
    \int_{q,r} (2\pi)^D \delta^{(D)}(p-q-l-r)\bigg[
      \delta^{(D)}(l)\cdot Y_{1,\nu}^{b}(p,q,r)
      \widetilde{B}_{\nu}^{b}(t,q)
      \\&+
    \frac{1}{2}
    Y_{2,\nu\rho}^{bc}(p,q,l,r)
    \widetilde{B}_{\nu}^{b}(t,q)
    \widetilde{B}_{\rho}^{c}(t,l)\bigg] \widetilde{\chi}(t,r)\,,
  \end{split}
\end{equation}
其中
\begin{equation}
  \begin{split}
  Y_{1,\nu}^b(q,r) &= \mathrm{i} \big( 2 r_\nu + (1-\kappa) q_\nu \big)
  T^b\,,\qquad
  Y_{2,\nu\rho}^{bc} = \delta_{\nu\rho} \bigl\{ T^{b}, T^{c} \bigr\}\,.
  \end{split}
\end{equation}
这些表达式给出流动夸克场的三点与四点顶点。对 $\bar{\chi}$ 可类似处理。


\subsection{费曼规则}
\label{ssec:Feynman_rules}


\subsubsection*{(流动的)传播子}

把动量空间中流动胶子场的解 \noeqn{eq:gluon_solution_momentum} 代入两点函数，
可得\,\cite{Luscher:2011bx}
\begin{align}
  \left\langle \widetilde{B}_\mu^a(t,p) \widetilde{B}_\nu^b(s,q)
  \right\rangle \Big|_{\textrm{LO}} &= \widetilde{K}_{\mu \rho}(t,p)
  \widetilde{K}_{\nu \sigma}(s,q) \left\langle \widetilde{A}_\rho^a(p)
  \widetilde{A}_\sigma^b(q) \right\rangle \nonumber \\ &= (2\pi)^D
  \delta^{(D)}(p+q) \, g^2\,D_{\mu\nu}^{ab}(p,t+s,\xi,\kappa)\,,
\end{align}
其中
\begin{equation}
  \begin{split}
    D_{\mu\nu}^{ab}(p,t,\xi,\kappa) =
  \delta^{ab} \frac{1}{p^2} \bigg( \Big(
  \delta_{\mu \nu} - \frac{p_\mu p_\nu}{p^2} \Big) e^{-tp^2} + \xi
  \frac{p_\mu p_\nu}{p^2} e^{-\kappa tp^2} \bigg)\,,
    %% \label{eq:}
  \end{split}
\end{equation}
且我们使用了基本胶子传播子的结果
\begin{equation}
\left\langle \widetilde{A}_\mu^a(p) \widetilde{A}_\nu^b(q) \right\rangle
\Big|_{\text{LO}} = (2\pi)^D \delta^{(D)}(p+q) \, g^2
D_{\mu\nu}^{ab}(p,0,\xi,0)\,.
\end{equation}
因子 $g^2$ 已计入下文相应的顶点之中。由于流动胶子传播子在 $t+s=0$ 时与基本
胶子传播子重合，我们可以用同一个费曼规则表示二者：
\begin{align}
  & \raisebox{-0.5em}{\includegraphics{images/gluon-propagator.pdf}}
  & = &
  \begin{split}
    D_{\mu\nu}^{ab}(p,t+s,\xi,\kappa) \,.
  \end{split}
  & &
  \label{eq:flowedgluon}
\end{align}
我们称之为（流动的）胶子传播子。\noeqn{eq:flowedgluon} 同样适用于混合传播子
$\langle \widetilde A \widetilde B \rangle$，只需令两个流时间变量之一为零。

对流动夸克场也可如法炮制。把其动量空间解 \noeqn{eq:quark_solution_momentum}
代入两点函数得到
\begin{align}
  \left. \left\langle \widetilde{\chi}^i_\alpha(t,p)
  \widetilde{\bar{\chi}}^j_\beta(s,q) \right\rangle
  \right|_{\text{LO}} &= \widetilde{K}(t,p) \widetilde{K}(s,q)
  \left\langle \widetilde{\psi}^i_\alpha(p)
  \widetilde{\bar{\psi}}^j_\beta(q) \right\rangle \nonumber \\ &=
  (2\pi)^D \delta^{(D)}(p+q)\, S^{ij}_{\text{F},\alpha\beta}(p,m,t+s)\,,
\end{align}
其中
\begin{equation}
  \begin{split}
S^{ij}_\mathrm{F}(p,m,t) &= \delta^{ij} \frac{-i\slashed{p}+m}{p^2+m^2}\,
e^{-tp^2},
  \end{split}
\end{equation}
且我们使用了基本夸克传播子的结果
\begin{equation}
  \left\langle \widetilde{\psi}^i_\alpha(p)
  \widetilde{\bar{\psi}}^j_\beta(q) \right\rangle = (2\pi)^D
  \delta^{(D)}(p+q)\,S_{\text{F},\alpha\beta}^{ij}(p,m,0)\,.
\end{equation}
既然可以把基本传播子表示成流时间为零时的流动传播子，二者便可用同一费曼规则
表示：
\begin{align}
  & \raisebox{-0.5em}{\includegraphics{images/quark-propagator.pdf}}
  & = &
  \begin{split}
    S_{\text{F},\alpha\beta}^{ij}(p,m,t+s) \,.
  \end{split}
  & &
\end{align}
我们称之为（流动的）夸克传播子，它同样涵盖混合传播子
$\langle\psi\bar\chi\rangle$ 与 $\langle\chi\bar\psi\rangle$。


\subsubsection*{流线}

由于拉格朗日乘子场 $L_\mu^a$、$\lambda$、$\bar{\lambda}$ 在
\noeqn{eq:Lagrangian_flow} 中没有二次型项，这些场没有传播子。然而，拉格朗日
量含有流动场与拉格朗日乘子场的双线性项。用标准方法可导出领先阶两点函数
\begin{equation}
  \left\langle B_\mu^a(t,x) L_\nu^b(s,y) \right\rangle = \delta^{ab}
  H_{\mu\nu}(t-s,x-y)\,,
\end{equation}
其中 $H_{\mu\nu}(t,x)$ 满足方程
\begin{equation}
  \left[(\partial_t - \partial_\sigma \partial_\sigma) \delta_{\mu\nu} + (1 -
  \kappa) \partial_\mu \partial_\nu\right] H_{\nu\rho}(t,x) =
  \frac{1}{2\TR} \delta_{\mu\rho} \delta(t) \delta^{(D)}(x)\,,
\end{equation}
$\TR$ 为通常的色迹归一化。由于基本胶子场 $A_\mu^a(x) = B_\mu^a(0,x)$ 不与拉格
朗日乘子场 $L_\mu^a$ 耦合，且所有流时间变量均为正，可以施加初始条件
\begin{equation}
  \begin{split}
  \left\langle B_\mu^a(t,x) L_\nu^b(s,y) \right\rangle
  \bigg|_{t=0} = 0 \quad \Rightarrow \quad
  H_{\mu\nu}(-s,x) = 0\,,
    \label{eq:bl}
  \end{split}
\end{equation}
于是唯一解为
\begin{equation}
  H_{\mu\nu}(t,x) = \frac{1}{2\TR} \theta(t) K_{\mu\nu}(t,x)\,.
  \label{eq:hmunu}
\end{equation}
我们把 $\langle BL\rangle$ 双线性称为``胶子流线''。因子 $1/(2\TR)$ 的倒数将
出现在下文相应的``流顶点''中，因此可将其完全弃去而写成
\begin{align}
  & \raisebox{-1.2em}{\includegraphics{images/gluon-flowline.pdf}}
  & = &
  \begin{split}
    \delta^{ab} \, \theta(t-s)\,\widetilde K_{\mu\nu}(t-s,p) \,,
  \end{split}
  & &
  \label{eq:blfr}
\end{align}
其中 $\widetilde K_{\mu\nu}(t,p)$ 已在 \noeqn{eq:kmunu} 定义，相邻箭头指示由
$\theta$ 函数所蕴含的指向增大流时间的方向。与真正的传播子不同，胶子流线对一
切 $p$ 都是正规函数。

类似地，混合费米子两点函数在领先阶确定为
\begin{equation}
  \left\langle \chi_\alpha^i(t,x) \bar{\lambda}_\beta^j(s,y) \right\rangle =
  \delta_{\alpha\beta} \delta^{ij} G(t-s,x-y)\,,
\end{equation}
其中 $G(t,x)$ 满足
\begin{equation}
  (\partial_t - \partial_\mu \partial_\mu) G(t,x) = \delta(t) \delta^{(D)}(x)\,,
\end{equation}
以及条件
\begin{equation}
  \left\langle \chi_\alpha^i(t,x) \bar{\lambda}_\beta^j(s,y)
  \right\rangle \bigg|_{t=0} = 0 \quad \Rightarrow \quad G(-s,x)= 0\,.
\end{equation}
唯一解为
\begin{equation}
  G(t,x) = \theta(t) K(t,x)\,.
\end{equation}
其傅里叶变换表达式定义了``费米子流线''费曼规则；$\theta$ 函数再次规定了方向，
如图中指向增大流时间的相邻箭头所示：
\begin{align}
  & \raisebox{-1.2em}{\includegraphics{images/quark-flowline.pdf}}
  & = &
  \begin{split}
    \delta_{\alpha\beta} \delta_{ij} \, \theta(t-s) \, \widetilde K(t-s,p) \,,
  \end{split}
  & &
  \label{eq:chibarlambda}
\end{align}
其中 $\widetilde K(t,p)$ 已在 \noeqn{eq:k} 定义。照例，费米子线\textit{上}的
箭头表示费米子的``电荷流''。这样我们对 $\langle\chi\bar\lambda\rangle$ 与
$\langle\lambda\bar\chi\rangle$ 两种双线性有了统一的费曼规则；后者可仿上得
到，恰好对应电荷流方向的反转：
\begin{align}
  & \raisebox{-1.2em}{\includegraphics{images/antiquark-flowline.pdf}}
  & = &
  \begin{split}
    \delta_{\alpha\beta} \delta_{ij} \, \theta(t-s) \, \widetilde K(t-s,p) \,.
  \end{split}
  & &
  \label{eq:lambdabarchi}
\end{align}
由于 $\widetilde K(t,p)$ 只通过 $p^2$ 依赖 $p$，费米子流线的动量方向无关紧要。


\subsubsection*{流顶点}

由于 \noeqn{eq:Lagrangian_flow} 中的 $\mathcal{L}_B$ 与 $\mathcal{L}_\chi$
正比于一个拉格朗日乘子场，所得顶点总至少涉及一条流线。这类顶点总是伴随一个
被积分的流时间参数。我们称之为``流顶点''，并在费曼图中用空心圆表示。相应的
费曼规则可直接推导。本文定义费曼规则时假定所有动量外向。

三点胶子流顶点由 $X_2$ 支配：
\begin{align}
  &
  \begin{gathered}
    \includegraphics{images/ggg-flow-vertex.pdf}
  \end{gathered}
  & = &
  \begin{split}
    -\mathrm{i} g f^{abc} \int_0^\infty \mathrm{d}s \, \big( & \delta_{\nu\rho} (r-q)_\mu + 2 \delta_{\mu\nu} q_\rho - 2 \delta_{\mu\rho} r_\nu \\
    & + (\kappa-1) (\delta_{\mu\rho} q_\nu  - \delta_{\mu\nu} r_\rho) \big) \,. \\
  \end{split}
  & &
  \label{eq:3gvert}
\end{align}
$\mathcal{L}_B$ 的相互作用项恰好含有一个拉格朗日乘子场 $L(s)$，参见
\noeqn{eq:Lagrangian_flow}。根据 \noeqn{eq:bl} 与 \noeqn{eq:blfr}，它必须与
更大流时间 $t>s$ 的流动胶子场 $B(t)$ 收缩。因此 \noeqn{eq:3gvert} 的顶点恰好
含一条外向流线。另一方面，$\mathcal{L}_B$ 相互作用项中的两个流动胶子场
$B(s)$ 的每一个既可以与\textit{更小}流时间 $t'<s$ 的拉格朗日乘子场 $L(t')$
收缩——形成内向流线——也可以与另一个流动胶子场 $B(s')$ 收缩。在此意义下，
费曼规则 \noeqn{eq:3gvert} 实际代表三个顶点，如 \fig{fig:3gcomb} 所示：它们
都含一条外向流线，而其余两条线中的每一条既可以是内向流线，也可以是流动胶子
传播子。\footnote{本文聚焦格林函数的计算；若计算振幅，这些线当然也可能代表
  外线``粒子''。}\noeqn{eq:3gvert} 中虚线箭头标示既可为流线亦可为传播子的
线，实心箭头则总表示流线。


\begin{figure}
  \begin{equation*}
    \begin{gathered}
      \includegraphics{images/ggg-flow-vertex_fff.pdf}
    \end{gathered}
    \qquad
    \begin{gathered}
      \includegraphics{images/ggg-flow-vertex_gff.pdf}
    \end{gathered}
    \qquad
    \begin{gathered}
      \includegraphics{images/ggg-flow-vertex_ggf.pdf}
    \end{gathered}
  \end{equation*}
  \caption{顶点 $X_2$ 的三个版本。带相邻箭头的线为流线，其余为流动的（或常规
    的）胶子。}
  \label{fig:3gcomb}
\end{figure}


类似地，四点胶子流顶点由 $X_3$ 支配：
\begin{align}
  &
  \begin{gathered}
    \includegraphics{images/gggg-flow-vertex.pdf}
  \end{gathered}
  & = &
  \begin{split}
    -g^2 \int_0^\infty \mathrm{d}s \, \big( & f^{abe} f^{cde} (\delta_{\mu\rho}\delta_{\nu\sigma} - \delta_{\mu\sigma}\delta_{\nu\rho} ) \\
    & + f^{ace} f^{bde} (\delta_{\mu\nu}\delta_{\rho\sigma} - \delta_{\mu\sigma}\delta_{\nu\rho} ) \\
    & + f^{ade} f^{bce} (\delta_{\mu\nu}\delta_{\rho\sigma} - \delta_{\mu\rho}\delta_{\nu\sigma} ) \big) \,.
  \end{split}
  & &
  \label{eq:4gvert}
\end{align}
同样，恰有一条外向流线，而其余三条线要么是流动胶子，要么是内向流线。这意味着
费曼规则 \noeqn{eq:4gvert} 实际代表四个不同的顶点。

注意，源自 \noeqn{eq:Lagrangian_flow} 中迹的因子 $2\TR$ 已按
\noeqn{eq:blfr} 的归一化在这两个顶点中被弃去。

将类似的考虑用于 $\mathcal{L}_\chi$ 可得到涉及流动夸克场的顶点。例如，带一
个胶子的夸克流顶点由如下费曼规则描述：
\begin{align}
  &
  \begin{gathered}
    \includegraphics{images/qqg-flow-vertex.pdf}
  \end{gathered}
  & = &
  \begin{split}
    \mathrm{i} g \, \delta_{\alpha\beta} \big( T^a \big)_{ij} \int_0^\infty \mathrm{d} s \, \big( 2 p_\mu + (1-\kappa) q_\mu \big) \,.
  \end{split}
  & &
\end{align}
此时，具有外向电荷流的费米子线在梯度流中同样是外向的，而另外两条线可以是流动
传播子或内向流线。于是该费曼规则实际代表四个不同的顶点。

完整的费曼规则列表见附录\,\ref{sec:frules}。

让我们总结 \gff\ 的费曼规则并指出若干特点：
\begin{enumerate}
  \item 对每个顶点，传播子与流线总是携带指数因子
    $\mathrm{e}^{\pm tp^2}$，其中 $t$ 为该顶点的流时间。
  \item 流线总始于某个流顶点，终于另一个流顶点或外部算符。$\theta$ 函数
    $\theta (t-s)$ 规定了流时间方向，我们在费曼图中以相邻箭头标示。
  \item 一个流顶点总有且仅有一条与之相连的外向流线。其余线或为内向流线，
    或为流动传播子。每个流顶点定义于流时间 $s$ 并蕴含积分 $\int_0^\infty
    \mathrm{d}s$；外向流线的 $\theta$ 函数把积分限制为有限上限 $t$，
    即 $\int_0^t \mathrm{d}s$。
  \item 含闭合流线圈的图为零，因为积分区间收缩为一个点。因此，只有流线构成
    树状结构的图才对任何可观测量有贡献。
\end{enumerate}
